% Fate's Edge SRD −−− Deck of Consequences (Optional) (Section 19)

\section{Deck of Consequences (Optional)}

The Deck of Consequences is a narrative tool for injecting drama, setbacks, and twists. It replaces or supplements GM fiat with randomized complications that remain thematically consistent. Its use is entirely optional.

\subsection{Deck Structure}

Use a standard 52−card deck (jokers optional). Suits map to complication themes:

\begin{tabular}{@{}llp{8cm}@{}}
⊤rule
\textbf{Suit} & \textbf{Theme} & \textbf{Examples} \\
∣rule
Hearts   & Social / Emotional & Betrayal, scandal, broken trust, awkward moment, ally under strain. \\
Spades   & Physical / Harm & Bruise, injury, collapse, gear damage, ambush. \\
Clubs    & Material / Cost & Resource loss, gear failure, debt, ruined supplies, theft. \\
Diamonds & Magical / Spiritual & Omen, curse, Patron attention, reality bend, haunting. \\
⊥tomrule
\end{tabular}

\subsection{Card Ranks & Severity}

\begin{itemize}
  ℹtem \textbf{Ace:} Scene−altering twist; compels immediate response.
  ℹtem \textbf{King, Queen, Jack:} Major complication with lasting effects.
  ℹtem \textbf{10–8:} Moderate complication that reshapes the current scene.
  ℹtem \textbf{7–5:} Minor complication; a nuisance that creates tension.
  ℹtem \textbf{4–2:} Subtle complication or foreshadowing omen.
\end{itemize}

\subsection{Jokers (Optional)}

\begin{itemize}
  ℹtem \textbf{Red Joker:} Catastrophic event (environmental collapse, Patron intervention).
  ℹtem \textbf{Black Joker:} Dark boon (immediate help, but with lasting cost or debt).
\end{itemize}

\subsection{When to Draw}

The GM may draw from the deck when:
\begin{itemize}
  ℹtem A roll shows multiple 1s (e.g., three or more).
  ℹtem The GM spends 2+ Story Beats at once.
  ℹtem A major timer fills (scene or campaign timer).
  ℹtem During travel to generate an unexpected encounter.
\end{itemize}

Aim for 1–2 draws per session; more if the tone skews chaotic.

\subsection{Procedure}

\begin{enumerate}
  ℹtem Note the number of SB spent (or generated).
  ℹtem Draw a number of cards equal to the SB spent (minimum 1, maximum 3).
  ℹtem Synthesize a single narrative complication from the combination of suits and the highest rank.
  ℹtem Do not reshuffle until the deck runs out or the session ends — this creates a natural ``fate budget.''
\end{enumerate}

⫕section{Synthesizing Complications}
Combine the suits and the highest rank into one fiction−first twist.

\textbf{Examples:}
\begin{itemize}
  ℹtem \textbf{Hearts + Spades:} Emotional betrayal that leads to physical harm (``Your ally flees, and the portcullis drops between you.'').
  ℹtem \textbf{Clubs + Diamonds:} Resource drain with magical backlash (``The ritual consumes your healing herbs and leaves a scar on the threshold.'').
  ℹtem \textbf{All four suits:} A crowd riots, loots your supplies, and a spirit feeds on the chaos.
\end{itemize}

\subsection{Shortcut Table for Fast Interpretation}

\begin{tabular}{@{}lll@{}}
⊤rule
\textbf{Suit} & \textbf{Low (2–5)} & \textbf{High (J–A)} \\
∣rule
Hearts   & Awkward social moment & Betrayal, scandal, broken trust \\
Spades   & Bruise, small stumble & Serious injury, structural collapse \\
Clubs    & Minor loss (rations, coins) & Major asset destroyed, debt called \\
Diamonds & Omen, eerie noise & Patron demand, reality bends \\
⊥tomrule
\end{tabular}

\subsection{Integration with Fiction}

\begin{itemize}
  ℹtem Translate the raw card result into a fiction−first complication. Do not apply mechanical penalties alone.
  ℹtem The complication should build on what is already happening in the scene, not restart the story.
  ℹtem Respect player agency — allow clever mitigation, but ensure the consequence lands.
\end{itemize}

\subsection{GM Quick Cues}

\begin{itemize}
  ℹtem Use the deck to surprise yourself as well as the players.
  ℹtem If a card does not fit the current fiction, draw again or adjust the rank up or down by one.
  ℹtem For campaign−scale foreshadowing, draw a card at the end of a session and note its theme; introduce the complication in a later session.
\end{itemize}

